
\section{Equilibrium suspended load in normal flow}

Water flows down a plane channel 300~m long and 10~m wide of slope
$S = 2\times10^{-4}$ with Manning roughness $n = 0.03$ at its normal depth
$h = 1$~m, so that gravity and friction balance and the flow is uniform,
$u = h^{2/3} S^{1/2}/n = 0.47$~m/s; Dirichlet boundaries at both ends hold
the normal state. Under the quadratic-drag shear closure with the Manning
friction factor, $\tau_b/\rho = f_c u^2$ with $f_c = g n^2 h^{-1/3}$, the
bed stress in normal flow is exactly $g h S$, so the Shields stress of the
100~$\mu$m sand is closed-form, $\tau^* = hS/(Rd) = 1.2$, thirty times the
critical value. That fixes the non-cohesive entrainment rate
$E = v_s\, 0.65\,\gamma_0 X/(1+\gamma_0 X)$, $X = \tau^*/\tau_c^* - 1$, and
against the well-mixed deposition $D = d^{*} v_s c$ the equilibrium
concentration $c_{eq} = E/(d^{*} v_s) = 0.043$. Clear water enters at the
inflow and the load relaxes to the equilibrium along the channel,
\begin{equation}
c(x) = c_{eq}\left(1 - e^{-x/L_s}\right), \qquad
L_s = \frac{q}{d^{*} v_s} = 59~\mathrm{m} .
\end{equation}
The bed is held fixed: the net exchange $E - D$ is largest at the inflow
and would otherwise scour the reach in minutes and take the depth, the
stress and the reference with it.

The case checks the quadratic-drag closure with the Manning friction
factor, the entrainment law and its balance with deposition in the
sediment mass balance under advection, that the flow stays the normal
flow, and that a fixed bed stays fixed. The steady profile is compared
cell by cell at the end of the run; the tracer advection is first-order
upwind, so the profile carries an error of order $\Delta x / L_s$.

\begin{figure}
\begin{center}
\includegraphics[width=0.9\textwidth]{concentration_plot.png}
\end{center}
\caption{Steady concentration along the channel, every cell, against
$c_{eq}(1 - e^{-x/L_s})$.}
\end{figure}

\begin{figure}
\begin{center}
\includegraphics[width=0.9\textwidth]{depth_plot.png}
\end{center}
\caption{Depth along the channel at the end of the run against the normal
depth.}
\end{figure}

\begin{figure}
\begin{center}
\includegraphics[width=0.9\textwidth]{approach_plot.png}
\end{center}
\caption{Approach to the steady state at four stations along the channel.}
\end{figure}

\endinput
